\section{Descent, wind response and suspension physics}
\label{sec:physics}

This section verifies Appendix F and the physical inputs to the wind
retrieval. Equations are consequences of the stated models, not proof that
an unmeasured parachute obeys those models. Numerical examples use proposal
values as design inputs, not flight measurements. Take $z$ positive upwards,
$s=z(0)-z$ positive downwards, $v_d=\dot s>0$, and
$g=9.80665\,\mathrm{m\,s^{-2}}$ for the numerical checks. Horizontal vehicle
velocity is $\boldsymbol u$, and horizontal air velocity is $\boldsymbol w$.
The suspension length $L_p$, turbulence scale $L$, and Obukhov length $L_O$
are different quantities.

\subsection{Terminal descent and canopy sizing}

For a deployed, approximately steady canopy, negligible buoyancy, no
vertical wind, and effective vertical drag area $D_z=C_{D,z}A$, Newton's
law in the downward direction is
\begin{equation}
 m\dot v_d=mg-\tfrac12\rho D_z v_d^2.
\end{equation}
Setting $\dot v_d=0$ gives
\begin{equation}
 v_t=\sqrt{\frac{2mg}{\rho D_z}},\qquad
 D_z=\frac{2mg}{\rho v_t^2},\qquad
 \rho=\frac{p}{R_dT_v}.
\end{equation}
Here $p$ is in pascals, $T_v$ is virtual temperature in kelvin, and
$R_d$ is the dry-air specific gas constant. In dry air $T_v=T$; neglecting
humidity otherwise is an approximation. $m$ must be the total mass supported
by the aerodynamic system, with the convention for canopy mass stated.
For a circular reference area $A=\pi d^2/4$,
\begin{equation}
 d=\sqrt{\frac{8mg}{\pi\rho C_{D,z}v_t^2}}.
\end{equation}
The area convention used with a quoted coefficient must match the convention
in this formula. Neither a numerical $C_D$ nor a finished canopy diameter is
established by the proposal.

If $v_{\min}\leq v_t\leq v_{\max}$ is a specified design band, then
\begin{equation}
 \frac{2mg}{\rho v_{\max}^2}\leq D_z
 \leq\frac{2mg}{\rho v_{\min}^2}.
\end{equation}
For uncertain $m,\rho,D_z$, the slowest terminal speed occurs at smallest
$m$ and largest $\rho D_z$, and the fastest at the opposite corner of the
uncertainty set. First-order fractional sensitivity is
\begin{equation}
 \frac{\delta v_t}{v_t}
 =\frac12\left(\frac{\delta m}{m}-\frac{\delta\rho}{\rho}
 -\frac{\delta C_D}{C_D}-\frac{\delta A}{A}\right).
\end{equation}
This is a differential identity; root-sum-square uncertainty requires
uncorrelated errors or the appropriate covariance terms.

\subsection{Transient and variable descent speed}

For constant parameters and release from rest, separating variables gives
\begin{align}
 \int_0^{v_d/v_t}\frac{dq}{1-q^2}&=\frac{gt}{v_t},\\
 v_d(t)&=v_t\tanh(gt/v_t),\\
 s(t)&=\frac{v_t^2}{g}\ln\cosh(gt/v_t).
\end{align}
Thus a $300\,\mathrm m$ descent at a constant
$10.5\,\mathrm{m\,s^{-1}}$ takes $28.57\,\mathrm s$, whereas the ideal
fully deployed release-from-rest model gives
\begin{equation}
 t(H)=\frac{v_t}{g}\operatorname{arcosh}
       \left[\exp\left(\frac{gH}{v_t^2}\right)\right]
 \simeq29.31\,\mathrm s\quad(H=300\,\mathrm m).
\end{equation}
The latter assumes instantaneous deployment and is not an inflation model.
Linearization around $v_t$ gives
\begin{equation}
 \delta\dot v=-\frac{2g}{v_t}\delta v,
 \qquad \tau_z=\frac{v_t}{2g}.
\end{equation}
The vertical relaxation time is half the horizontal time in the isotropic
drag approximation derived below.

For variable density or drag, integrate the dynamical equation with
$\dot z=-v_d$ instead of imposing instantaneous terminal equilibrium.
If quasi-steady descent is justified, then
\begin{equation}
 t_{\rm remaining}=\int_0^h\frac{dz}{v_d(z)},\qquad
 \Delta\boldsymbol x=\int_0^h\frac{\boldsymbol u(z)}{v_d(z)}\,dz.
\end{equation}
Quasi-steady tracking requires the fractional equilibrium-speed change
over $\tau_z$ to be small. A vertical air velocity $w_z$, positive upwards,
changes the air-relative downward speed to $v_d+w_z$; ground-relative
descent speed cannot then be substituted indiscriminately in drag.
Perturbing the first integral yields
\begin{equation}
 \delta t=\frac{\delta h}{v_d(h)}
 -\int_0^h\frac{\delta v_d(z)}{v_d(z)^2}\,dz.
\end{equation}
Altitude and descent-model uncertainty therefore both enter touchdown time.

\subsection{Pressure altitude: the precise hypsometric equation}

Hydrostatic balance $dp/dz=-\rho g$ and the ideal gas relation give
\begin{equation}
 dz=-\frac{R_dT_v}{g}\,d\ln p.
\end{equation}
With approximately constant $g$ over the layer,
\begin{align}
 z_2-z_1&=\frac{R_d}{g}\int_{p_2}^{p_1}T_v(p)\,d\ln p\\
 &=\frac{R_d\overline T_v^{\ln p}}{g}\ln(p_1/p_2),\\
 \overline T_v^{\ln p}&=
 \frac{\int_{p_2}^{p_1}T_v(p)\,d\ln p}{\ln(p_1/p_2)}.
\end{align}
The exact mean is weighted in log pressure, rather than an unspecified
arithmetic mean. In the isothermal case,
$p(z)=p_1\exp[-g(z-z_1)/(R_dT_v)]$; substitution recovers $z-z_1$
exactly. Differentiation of the constant-temperature form gives
\begin{equation}
 \delta h=\frac{R_d}{g}\left[
  \ln(p_1/p_2)\delta T_v+
  T_v\left(\frac{\delta p_1}{p_1}-\frac{\delta p_2}{p_2}\right)\right].
\end{equation}
A common pressure bias may largely cancel in a pressure difference; treating
all samples as independent can therefore give the wrong uncertainty.
For air density, the dry-air differential is
\begin{equation}
 \frac{\delta\rho}{\rho}=\frac{\delta p}{p}-\frac{\delta T_v}{T_v},
 \qquad
 \frac{\operatorname{Var}(\rho)}{\rho^2}\simeq
 \frac{\operatorname{Var}(p)}{p^2}
 +\frac{\operatorname{Var}(T_v)}{T_v^2}
 -\frac{2\operatorname{Cov}(p,T_v)}{pT_v}.
\end{equation}
Humidity uncertainty enters through $T_v$. In the same ideal-gas,
hydrostatic model,
\begin{equation}
 \frac{d\ln\rho}{dz}=-\frac{g}{R_dT_v}
                         -\frac{1}{T_v}\frac{dT_v}{dz}.
\end{equation}
An uncertain assigned height consequently perturbs a prescribed density
profile by $\delta\rho\simeq\rho_z\delta z$. Do not add that term again
if the density is already estimated directly from the same uncertain
pressure--temperature observations in a joint model: their covariance
must be propagated together.
Pressure altitude needs a reference pressure and a ground elevation to
become altitude above ground. Dynamic pressure at a sensor inlet, thermal
offsets, and non-hydrostatic motions are separate errors, not removed by
the algebra.

\subsection{Horizontal lag: what follows from isotropic drag}

Let $\boldsymbol q=(\boldsymbol u-\boldsymbol w,-v_d)$ denote air-relative
velocity with zero vertical wind. Suppose the total aerodynamic force is
the isotropic vector model
\begin{equation}
 \boldsymbol F_a=-b\lvert\boldsymbol q\rvert\boldsymbol q,
 \qquad b=\tfrac12\rho C_DA.
\end{equation}
For $\lvert\boldsymbol u-\boldsymbol w\rvert\ll v_d$,
$\lvert\boldsymbol q\rvert=v_d+O(\lvert\boldsymbol u-\boldsymbol w\rvert^2/v_d)$.
The horizontal equation is consequently
\begin{equation}
 m\dot{\boldsymbol u}=-bv_d(\boldsymbol u-\boldsymbol w).
\end{equation}
At vertical terminal equilibrium $bv_d^2=mg$, so
\begin{equation}
 \tau_h=\frac{m}{bv_d}=\frac{v_d}{g},\qquad
 \ell_h=v_d\tau_h=\frac{v_d^2}{g}.
\end{equation}
At $v_d=10.5\,\mathrm{m\,s^{-1}}$ this gives
$\tau_h=1.07070\,\mathrm s$ and $\ell_h=11.2424\,\mathrm m$.
Because $d/dt=v_d\,d/ds$, the first-order model implies
\begin{equation}
 \boldsymbol w(s)=\boldsymbol u(s)+\ell_h(s)
                       \frac{d\boldsymbol u}{ds}.
\end{equation}
There is no derivative of $\ell_h$ in this algebraic inversion. A
time-dependent coefficient does prevent the use of a single global Fourier
transfer function.

\subsection{Why canopy terminal speed does not identify horizontal lag}

A real canopy has an orientation and different axial, normal, and rotational
responses. Define its local horizontal force Jacobian by
\begin{equation}
 \boldsymbol F_{a,h}\simeq
 -K_h(\boldsymbol u-\boldsymbol w)+\boldsymbol F_0,
 \qquad
 M_h\dot{\boldsymbol u}
 =-K_h(\boldsymbol u-\boldsymbol w)+\boldsymbol F_0.
\end{equation}
Here $M_h$ may include an effective added-mass contribution when justified,
and $\boldsymbol F_0$ represents trim or coupling forces. Provided $K_h$ is
invertible, the conditional local inverse is
\begin{equation}
 \boldsymbol w=\boldsymbol u+K_h^{-1}M_h\dot{\boldsymbol u}
                            -K_h^{-1}\boldsymbol F_0.
\end{equation}
Neither $K_h$ nor $\boldsymbol F_0$ follows from $F_{a,z}=mg$.
As an explicit counterexample, parameterize one axis by
$K_h=rmg/v_d$, $M_h=m$, $\boldsymbol F_0=0$, with $r>0$.
Every value of $r$ is consistent with the same prescribed vertical
terminal balance, while
\begin{equation}
 \tau_h=\frac{v_d}{rg},\qquad \ell_h=\frac{v_d^2}{rg}.
\end{equation}
The illustrative ratios $r=0.5,1,2$ give relaxation times
$2.14140,1.07070,0.535351\,\mathrm s$. These are sensitivity cases, not
measured coefficients. Thus Appendix F.3 is a valid derivation within its
isotropic model, but the claimed physical lag length is not established for
the chosen canopy. Identification from instrumented drops is needed.

Primary-source context: the NASA study by Pei et al., \emph{Modal Analysis
of the Orion Capsule Two Parachute System}, uses separate axial and normal
force coefficients, entrapped-air mass, and canopy-dependent oscillation
parameters. Its multi-canopy configuration is not this CanSat; it demonstrates
why a single drag coefficient cannot establish all lateral dynamics.
\url{https://ntrs.nasa.gov/api/citations/20200002721/downloads/20200002721.pdf}

\subsection{Frequency response, resolution and pendulum contamination}

For constant $\ell_h$, a wind Fourier component $w(s)=W e^{iks}$ gives
\begin{equation}
 \frac{U}{W}=\frac{1}{1+ik\ell_h},\qquad
 \lvert H^{-1}\rvert=\sqrt{1+(k\ell_h)^2},\quad k=2\pi/\lambda.
\end{equation}
This verifies the inversion-amplification factor. At wavelength
$\lambda=20\,\mathrm m$ and $\ell_h=11.2424\,\mathrm m$, amplification
is approximately $3.670$. A $20\,\mathrm m$ bin is not automatically a
resolved $20\,\mathrm m$ Fourier structure. Resolution must be read from
the retrieval's averaging kernels and error covariance.

If suspension motion adds velocity $u_p$ to the canopy or system-reference
velocity $u_c$, then applying the same inverse to $u=u_c+u_p$ adds
the false wind term $u_p+\tau_h\dot u_p$.
The required correction is a translational suspension or lever-arm velocity,
not merely a camera rotation correction. For a measured rigid displacement
$\boldsymbol r=R\boldsymbol r_b$ from the reference point,
\begin{equation}
 \boldsymbol v_{\rm sensor}=\boldsymbol v_{\mathrm{ref}}
                  +R(\boldsymbol\omega_b\times\boldsymbol r_b).
\end{equation}
A flexible suspension needs its own configuration estimate. GPS and image
sensors can share this physical bias; agreement between their inferred winds
does not by itself validate absolute wind accuracy.

\subsection{Pendulum specific force and the limits of gyro acceleration}

For a point payload, let $\boldsymbol e$ point from a moving pivot to the
payload, $\boldsymbol r=\boldsymbol r_P+L_p\boldsymbol e$, and let tension
be $-T\boldsymbol e$. With gravity $\boldsymbol g$ and other force
$\boldsymbol F_o$,
\begin{equation}
 \boldsymbol a=\boldsymbol g-\frac{T}{m}\boldsymbol e
                         +\frac{\boldsymbol F_o}{m},\qquad
 \boldsymbol f=\boldsymbol a-\boldsymbol g
             =-\frac{T}{m}\boldsymbol e+\frac{\boldsymbol F_o}{m}.
\end{equation}
Thus the ideal accelerometer measures tension direction, rather than gravity
direction. Lateral \emph{body-frame} readings vanish only when the measured
body axis is aligned with the suspension, the sensor is at the modeled
point, and other forces vanish. Body rotation relative to the line and
rotational lever-arm accelerations violate that special case. Appendix
F.4's statement that residuals are necessarily too small to help is not
supported by a measurement or model.

In a planar model with $x$ horizontal and $z$ up,
$\boldsymbol e=(\sin\theta,-\cos\theta)$ and
$\boldsymbol e_\theta=(\cos\theta,\sin\theta)$.
Projecting acceleration onto $\boldsymbol e_\theta$ eliminates tension:
\begin{equation}
 L_p\ddot\theta=-g\sin\theta
 -a_{P,x}\cos\theta-a_{P,z}\sin\theta
 +\frac{\boldsymbol F_o\cdot\boldsymbol e_\theta}{m}.
\end{equation}
After small-angle linearization and an assumed viscous damping law,
\begin{equation}
 \ddot\theta+2\zeta_p\omega_0\dot\theta+\omega_0^2\theta=F(t),
 \qquad \omega_0^2=\frac{g+a_{P,z}}{L_p},
\end{equation}
when $a_{P,z}$ is sufficiently constant and the remaining forcing is
collected into $F$, with units $\mathrm{s^{-2}}$.
$\omega_0^2=g/L_p$ is the fixed-pivot limit. Aerodynamic canopy motion can
change both effective pivot and restoring dynamics.

Rearrangement gives the exact linear-model diagnostic
\begin{equation}
 \theta=\frac{F-\ddot\theta-2\zeta_p\omega_0\dot\theta}{\omega_0^2}.
\end{equation}
The proposed $\theta\simeq-\ddot\theta/\omega_0^2$ follows only when
forcing and damping contributions are negligible. Constant forcing admits
$\theta=F/\omega_0^2$ with zero angular velocity and acceleration: a
direct counterexample to unconditional tilt observability from gyro
acceleration. Differentiating a noisy gyro also amplifies high-frequency
noise. This can be a calibrated model diagnostic, not an independent
absolute attitude measurement without additional assumptions.

For a stable, constant-coefficient linear oscillator with stationary
forcing and transients removed, Fourier transformation gives
\begin{equation}
 H_p(\omega)=\frac{1}{\omega_0^2-\omega^2+2i\zeta_p\omega_0\omega},
 \qquad S_\theta(\omega)=|H_p(\omega)|^2S_F(\omega).
\end{equation}
This verifies F.5's spectrum relation under those conditions. Measured
swing growth alone cannot distinguish changing forcing from changing
canopy damping or natural frequency.

\subsection{Shear, the sampling path, and a conditional frozen field}

For perfect local mean-wind tracking, $U(z)=U_0+Sz$ and
$z(t)=z_0-v_dt$ imply
\begin{align}
 x(t)-x(0)&=\int_0^tU(z(t'))\,dt',\\
 \xi(t)&=x(t)-x(0)-tU(z(t))\\
 &=\int_0^t[U(z(t'))-U(z(t))]\,dt'
 =\frac{Sv_dt^2}{2}=\frac{S(\Delta z)^2}{2v_d}.
\end{align}
The subtraction compares the trajectory to uniform advection at the
current layer's speed. It is not a single global frame comoving with
every sheared layer. For $S=0.05\,\mathrm{s^{-1}}$, $v_d=10.5\,\mathrm{m\,s^{-1}}$
and $\Delta z=100\,\mathrm m$, $\xi=23.81\,\mathrm m$, or $23.8\%$
of the vertical separation. At $\Delta z=10\,\mathrm m$, $\xi=0.2381\,\mathrm m$,
or $2.38\%$. F.6's arithmetic is correct at $100\,\mathrm m$, but its
``few percent'' conclusion fails at that separation.

With finite lag, $\tau_h\dot u+u=U(t)$ and
$U(t)=U(0)-Sv_dt$ have solution
\begin{equation}
 u(t)=U(t)+\tau_hSv_d+
 [u(0)-U(0)-\tau_hSv_d]e^{-t/\tau_h}.
\end{equation}
The locally air-relative path slope after the transient is
$|u-U|/v_d=|S|\tau_h$, which is $0.05354$ in the isotropic numerical
example. It is distinct from the finite-separation offset above.
An approximately vertical covariance model needs both effects, suspension
motion, turbulence and vertical wind to be small on the separations used.

If an eddy has crossing scale $L$, characteristic internal velocity
$\sigma$, and turnover estimate $t_e\sim L/\sigma$, then
\begin{equation}
 \frac{t_{\mathrm{cross}}}{t_e}\sim\frac{\sigma}{v_d},\qquad
 |S|t_{\mathrm{cross}}\sim\frac{|S|L}{v_d}.
\end{equation}
The first estimate being small is helpful, but not sufficient: shear
deformation, anisotropic eddy scales, scale-dependent sweeping, evolution,
nonstationarity and nonvertical sampling also matter. The ratio is
scale-independent only if the same $L$ and a scale-independent $\sigma$
describe both estimates. Values $0.05$--$0.1$ are hypotheses until measured.
F.7's claim of validity at every height is therefore not proved.
Direct atmospheric measurements and simulations have found failures of the
frozen-field assumption at some scales: Cheng et al. (2017),
\url{https://doi.org/10.1002/2017GL073499}.

\section{Thermodynamic stability and surface-layer inference}
\label{sec:boundary}

\subsection{Potential temperature from the first law}

For a dry ideal gas, $h=c_pT$ and specific volume $\alpha=R_dT/p$.
Reversible adiabatic motion satisfies $dh-\alpha\,dp=0$, hence
\begin{equation}
 c_p\,dT-R_dT\frac{dp}{p}=0,
 \qquad d\ln T=\frac{R_d}{c_p}d\ln p.
\end{equation}
Integration defines the conserved potential temperature
\begin{equation}
 \theta=T\left(\frac{p_0}{p}\right)^{R_d/c_p},
 \qquad p_0=1000\,\mathrm{hPa},\quad R_d/c_p\simeq0.286.
\end{equation}
The pressure ratio must use consistent units. Under hydrostatic balance,
\begin{equation}
 \frac{d\ln\theta}{dz}
 =\frac{1}{T}\frac{dT}{dz}+\frac{g}{c_pT},
 \qquad
 \frac{d\theta}{dz}=\frac{\theta}{T}
                    \left(\frac{dT}{dz}+\frac{g}{c_p}\right).
\end{equation}
For $c_p=1004\,\mathrm{J\,kg^{-1}K^{-1}}$, the dry-adiabatic lapse
is $g/c_p=9.7676\,\mathrm{K\,km^{-1}}$, or $2.9303\,\mathrm K$
over $300\,\mathrm m$. A temperature decrease with altitude can still
have increasing $\theta$, and therefore stable dry stratification.

A parcel displaced upwards by $\delta z$ approximately preserves its
$\theta$ while ambient $\theta$ changes by $\theta_z\delta z$.
Its buoyancy is $-(g/\theta)\theta_z\delta z$, so
\begin{equation}
 \delta\ddot z=-N^2\delta z,\qquad
 N^2=\frac{g}{\theta}\frac{d\theta}{dz}.
\end{equation}
$N^2>0$ gives restoring buoyancy; $N^2<0$ gives local dry static
instability. Unsaturated moist-air buoyancy uses virtual potential
temperature; saturated parcel stability needs moist thermodynamics.
The dry expressions do not settle those cases.

Temperature and pressure uncertainty propagate as
\begin{equation}
 \frac{\delta\theta}{\theta}
 =\frac{\delta T}{T}-\frac{R_d}{c_p}\frac{\delta p}{p},
 \qquad \operatorname{Var}(\theta)\simeq
 \boldsymbol J_\theta\Sigma_{T,p}\boldsymbol J_\theta^{\mathsf T}.
\end{equation}
Differences between levels require their cross-covariances. A first-order
thermistor model obeys
\begin{equation}
 \tau_T\dot T_m+T_m=T_{\rm air},\qquad
 T_{\rm air}=T_m+v_d\tau_T\frac{dT_m}{ds}.
\end{equation}
Its spatial lag length is $v_d\tau_T$, equal to $10.5\,\mathrm m$ if
$\tau_T=1\,\mathrm s$ and the proposal speed applies. Differentiation
again needs regularization; shielding does not determine response time or
eliminate all radiation and self-heating error.

\subsection{MOST: integration with displacement and the lower boundary}

Monin--Obukhov similarity theory (MOST) describes ensemble or suitable
time-averaged profiles over approximately homogeneous terrain in a
stationary, constant-flux surface layer. Let $d_0$ be displacement height,
$h_e=z-d_0>0$, $z_0$ aerodynamic roughness length, and
\begin{equation}
 u_*=[(\overline{u'w'})^2+(\overline{v'w'})^2]^{1/4},\qquad
 L_O=-\frac{u_*^3\overline\theta_v}
            {\kappa g\overline{w'\theta_v'}}.
\end{equation}
For an aligned mean wind $U$, the assumed flux-gradient relationship is
\begin{equation}
 \frac{dU}{dz}=\frac{u_*}{\kappa h_e}\phi_m(h_e/L_O).
\end{equation}
Define $\psi_m(0)=0$ and
\begin{equation}
 \psi_m'(\zeta)=\frac{1-\phi_m(\zeta)}{\zeta}.
\end{equation}
Integrating from the aerodynamic reference $h_e=z_0$, where the
extrapolated profile has $U=0$, yields
\begin{equation}
 U(z)=\frac{u_*}{\kappa}
 \left[\ln\frac{z-d_0}{z_0}-\psi_m\left(\frac{z-d_0}{L_O}\right)
                         +\psi_m\left(\frac{z_0}{L_O}\right)\right].
\end{equation}
The lower-boundary term and displacement height are absent from K.2.
Dropping the former requires $|z_0/L_O|\ll1$; dropping the latter
requires justification from terrain. The zero condition is an extrapolated
aerodynamic definition, not a literal measurement at a plant surface.
Avoid the logarithmic singularity and avoid extrapolating MOST to touchdown.
A reference-height difference can be fitted without specifying zero wind
at ground:
\begin{equation}
 U(z_2)-U(z_1)=\frac{u_*}{\kappa}
 \left[\ln\frac{z_2-d_0}{z_1-d_0}
 -\psi_m\left(\frac{z_2-d_0}{L_O}\right)
 +\psi_m\left(\frac{z_1-d_0}{L_O}\right)\right].
\end{equation}

\subsection{Checking the empirical stability corrections}

The frequently used Dyer choices are
\begin{equation}
 \phi_m(\zeta)=
 \begin{cases}
  1+5\zeta,&\zeta\geq0,\\
  (1-16\zeta)^{-1/4},&\zeta<0.
 \end{cases}
\end{equation}
The constants are empirical, not first-principles consequences of MOST.
For the stable branch integration gives $\psi_m=-5\zeta$.
For the unstable branch put $x=(1-16\zeta)^{1/4}$; then
$d\zeta=-x^3dx/4$. Substitution into $d\psi_m=(1-1/x)d\zeta/\zeta$
gives
\begin{equation}
 d\psi_m=\frac{4x^2}{(x+1)(x^2+1)}\,dx
 =\left(\frac{2}{x+1}+\frac{2x-2}{x^2+1}\right)dx.
\end{equation}
Integrating from $x=1$ enforces $\psi_m(0)=0$:
\begin{equation}
 \psi_m=2\ln\frac{1+x}{2}+\ln\frac{1+x^2}{2}
        -2\arctan x+\frac\pi2.
\end{equation}
Differentiation recovers $1-\zeta\psi_m'=\phi_m$, checking the
normalization and signs. In the neutral limit $L_O\to\pm\infty$,
$\psi_m\to0$ and the logarithmic law is recovered. Positive heat flux
gives negative $L_O$ under the convention above.

The original field experiment is Businger et al. (1971),
\url{https://journals.ametsoc.org/view/journals/atsc/28/2/1520-0469_1971_028_0181_fprita_2_0_co_2.xml}.
The specific $5$ and $16$ choices and their empirical status are discussed
by Dyer (1974), primary paper available at
\url{https://gibbs.science/efd/handouts/dyer_1974.pdf}.
No fixed $50$--$100\,\mathrm m$ layer is guaranteed to meet the
constant-flux, homogeneous-surface assumptions on launch day. A single
short descent through different air parcels is not automatically a mean
wind profile: drift, temporal variability and heterogeneous fetch can
confound the fit.

\subsection{Richardson numbers and what the quarter threshold means}

Comparing stratification with vector shear gives
\begin{equation}
 Ri_g=\frac{N^2}{S^2},\qquad
 S^2=\left(\frac{d\overline u}{dz}\right)^2
       +\left(\frac{d\overline v}{dz}\right)^2.
\end{equation}
Approximating both derivatives across a layer of thickness $\Delta z>0$
gives the proposal's bulk quantity
\begin{equation}
 Ri_b=\frac{g}{\overline\theta}
 \frac{\Delta\theta\,\Delta z}
      {(\Delta\overline u)^2+(\Delta\overline v)^2}.
\end{equation}
It is dimensionless. Negative $\Delta\theta$ diagnoses dry static
instability under the earlier assumptions. A nearly zero denominator
makes this ratio poorly determined; report shear and stratification
uncertainties rather than forcing a classification. Turning wind requires
both velocity components, not speed differences alone.

For explicit uncertainty propagation, set $A=g/\overline\theta$,
$q=\Delta\theta$, $h=\Delta z$, $a=\Delta\overline u$,
$b=\Delta\overline v$, and $D=a^2+b^2$. With $r=Ri_b=Aqh/D$, the
nonzero partial derivatives are
\begin{align}
 r_q&=Ah/D,& r_h&=Aq/D,\\
 r_a&=-2Aqha/D^2,&r_b&=-2Aqhb/D^2,\\
 r_{\overline\theta}&=-r/\overline\theta.
\end{align}
Thus $\operatorname{Var}(r)\simeq\boldsymbol J_r\Sigma
\boldsymbol J_r^{\mathsf T}$ for the joint level-difference, mean-temperature
and height covariance. This expression remains well behaved at $q=0$
when $D>0$, whereas fractional-error formulas do not. Near $D=0$, even
this linearization can fail and a joint Monte Carlo or interval treatment
is required. Adjacent layers share sensor samples, pressure-derived
heights, thermistor lag and retrieved wind errors, so neither their
$Ri_b$ values nor the errors in numerator and denominator are generally
independent.

The Miles--Howard theorem is a statement about infinitesimal normal-mode
stability of steady, parallel, inviscid, nondiffusive, stably stratified
flow. The sufficient condition is $Ri_g>1/4$ everywhere, not a condition
on a coarse $Ri_b$, and not a guarantee that existing turbulence vanishes.
A compact independent algebra check follows. For the two-dimensional
Taylor--Goldstein problem with perturbation $\varphi(z)e^{ik(x-ct)}$,
\begin{equation}
 \varphi''+\left[\frac{N^2}{(U-c)^2}
 -\frac{U''}{U-c}-k^2\right]\varphi=0.
\end{equation}
Suppose $c_i=\operatorname{Im}c>0$, so $U-c$ has no real zero, and set
$\varphi=(U-c)^{1/2}F$. Differentiation gives
\begin{equation}
 [(U-c)F']'+\left[\frac{N^2-U'^2/4}{U-c}
              -\frac{U''}{2}-k^2(U-c)\right]F=0.
\end{equation}
Multiply by $F^*$, integrate with homogeneous boundary values (or
adequate decay), and take the imaginary part. The result is
\begin{equation}
 c_i\int\left[|F'|^2+k^2|F|^2+
       \frac{N^2-U'^2/4}{|U-c|^2}|F|^2\right]dz=0.
\end{equation}
For a nonzero mode with $k>0$, $N^2-U'^2/4\geq0$ makes the integral
positive, contradicting $c_i>0$. Conversely, violating this sufficient
condition does not prove instability. Finite amplitude, viscosity,
diffusion, forcing, time dependence and coarse averaging require separate
analysis. Howard's original \emph{Note on a paper of John W. Miles},
\emph{Journal of Fluid Mechanics} 10(4), 509--512 (1961), is
\url{https://doi.org/10.1017/S0022112061000317}; the conditions and distinctions
are also stated in the primary research of Liu, Thorpe and Smyth (2012),
\url{https://www.pmf.unizg.hr/_download/repository/Liu_Thorpe_Smyth_turb-stratif-shearJFM2012.pdf}.
Consequently K.3's claim that bulk values above $0.25$ establish insufficient
shear to sustain turbulence must be weakened to a conditional stability
diagnostic.

\subsection{Why a temperature lapse does not determine Obukhov length}

$L_O$ contains momentum and buoyancy \emph{fluxes}; a temperature gradient
alone does not measure either. A model-dependent relation can be derived
if MOST applies and both gradients are measured. With
$\theta_*=-\overline{w'\theta'}/u_*$ and
\begin{equation}
 \frac{d\theta}{dz}=\frac{\theta_*}{\kappa h_e}\phi_h(\zeta),
 \qquad \zeta=h_e/L_O,
\end{equation}
substitution into $Ri_g$ for aligned shear gives
\begin{equation}
 Ri_g=\zeta\frac{\phi_h(\zeta)}{\phi_m(\zeta)^2}.
\end{equation}
For the assumed stable Dyer choice $\phi_h=\phi_m=1+5\zeta$,
\begin{equation}
 Ri_g=\frac{\zeta}{1+5\zeta},\qquad
 \zeta=\frac{Ri_g}{1-5Ri_g}\quad(0\leq Ri_g<0.2).
\end{equation}
The $0.2$ limit is a property of these closures; it is different from the
Miles--Howard quarter threshold. Inversion is ill-conditioned near it.
Replacing gradient $Ri_g$ with bulk $Ri_b$ without integrated profile
relationships is not an exact derivation. S2 can test consistency between
joint wind--temperature profiles and a fitted stability model; it cannot
claim that temperature lapse alone fixes the relevant MOST class or $L_O$.

\subsection{Neutral profile fitting and roughness uncertainty}

For neutral conditions, specified $d_0$, and a fixed reference length
$h_r$ to keep logarithms dimensionless,
\begin{equation}
 U(z)=\alpha\ln\frac{z-d_0}{h_r}+\beta,
 \qquad u_*=\kappa\alpha,\qquad
 z_0=h_r\exp(-\beta/\alpha).
\end{equation}
For observations $\boldsymbol y=X\boldsymbol b+\boldsymbol\epsilon$,
$\boldsymbol b=(\alpha,\beta)^{\mathsf T}$ and known covariance $\Sigma$,
minimizing the weighted residual quadratic gives
\begin{equation}
 \widehat{\boldsymbol b}=(X^{\mathsf T}\Sigma^{-1}X)^{-1}
                   X^{\mathsf T}\Sigma^{-1}\boldsymbol y,
 \quad C_b=(X^{\mathsf T}\Sigma^{-1}X)^{-1}.
\end{equation}
The formula requires rank two and assumes heights and displacement are
known; uncertain heights require errors-in-variables or a joint model.
With $q=\ln(z_0/h_r)=-\beta/\alpha$,
\begin{equation}
 \nabla q=\left(\frac{\beta}{\alpha^2},-\frac1\alpha\right),
 \qquad \operatorname{Var}(q)\simeq\nabla q\,C_b\,\nabla q^{\mathsf T}.
\end{equation}
A near-Gaussian $q$ produces multiplicative, asymmetric intervals for
$z_0$. Small $\alpha$, limited height range, unknown $d_0$, changing
stability and correlated profile errors reduce identifiability. The
land-cover values in K.5 can be broad priors; they are not site-specific
measurements and cannot establish a validated roughness estimate.

\section{Supporting engineering arithmetic}
\label{sec:engineering}

\subsection{Footprint and exposure blur}

Assume a rectilinear nadir camera over horizontal terrain, horizontal
field of view $90^\circ$, square pixels and $1280\times800$ resolution.
Then $f_x=f_y=640$ pixels, normalized horizontal image edge is $|x|=1$,
and vertical edge is $|y|=0.625$. The ground width and height are
\begin{equation}
 W_g=2h\tan45^\circ=2h,\qquad
 H_g=\frac{800}{1280}W_g=1.25h,
 \quad \mathrm{GSD}=h/f=0.46875\,\mathrm{m/px}\quad(h=300\,\mathrm m).
\end{equation}
Thus the quoted $600\times375\,\mathrm m$ footprint is consistent with
a \emph{horizontal}, rather than diagonal, $90^\circ$ field. A real lens
needs distortion calibration and the actual field specification.

For a stationary scene point with camera coordinates $(X,Y,Z)$,
$x=X/Z$, $y=Y/Z$, camera translation $\boldsymbol v$ and angular velocity
$\boldsymbol\omega$, $\dot{\boldsymbol P}=-\boldsymbol v-
\boldsymbol\omega\times\boldsymbol P$ gives
\begin{align}
 \dot x&=\frac{-v_x+xv_z}{Z}+xy\omega_x-(1+x^2)\omega_y+y\omega_z,\\
 \dot y&=\frac{-v_y+yv_z}{Z}+(1+y^2)\omega_x-xy\omega_y-x\omega_z.
\end{align}
For exposure $t_e$ short enough that velocity is approximately constant,
the path length in pixels is
\begin{equation}
 B\simeq t_e\sqrt{(f_x\dot x)^2+(f_y\dot y)^2}.
\end{equation}
This is a flow approximation; the actual image is the time integral of the
moving projection. Translation, pitch, roll, spin and descent contribute
simultaneously, and their vectors may add or cancel.

At a pitch rate $40^\circ/\mathrm s=0.698132\,\mathrm{rad/s}$ and
$t_e=1\,\mathrm{ms}$, pure pitch gives $0.4468$ pixels at image center,
$0.8936$ pixels at the horizontal edge, and $0.9362$ pixels at the corner.
Pure spin at one revolution per second gives
\begin{equation}
 B=r_{\rm px}(2\pi)t_e,
 \qquad r_{\rm px}=\sqrt{640^2+400^2},\qquad B=4.7420\,\mathrm{px}
\end{equation}
at that corner, halved for a $0.5\,\mathrm{ms}$ exposure.
Descent adds radial flow of magnitude approximately
$r_{\rm px}v_d/h$: at $h=80\,\mathrm m$, $v_d=10.5\,\mathrm{m/s}$
and $1\,\mathrm{ms}$ it contributes $0.0991$ pixels at the corner.
The proposal's subpixel swing claim is consistent with isolated pitch
under these ideal intrinsics, but is not a bound on combined motion.
Short exposure also requires adequate photon signal; daylight alone does
not verify signal-to-noise ratio or shutter timing.

\subsection{Frame count, storage and processing capacity}

With $N_x,N_y$ pixels, stored $b$ bits per pixel and frame rate $f_c$,
\begin{equation}
 B_{\rm frame}=N_xN_yb/8,
 \quad R_{\rm image}=f_cB_{\rm frame},
 \quad S_{\rm total}=N_fB_{\rm frame}+S_{\rm metadata}+S_{\rm sensors}.
\end{equation}
Compression changes these numbers only if its achieved size is measured.
At $1280\times800$, eight-bit grayscale and $5\,\mathrm{Hz}$,
$B_{\rm frame}=1{,}024{,}000$ bytes and $R_{\rm image}=5.12$ decimal
MB/s. A $30\,\mathrm s$ image record is $153.6$ MB; a $60\,\mathrm s$
record is $307.2$ MB. The proposal's $300$ MB per flight therefore
needs a longer logging window, sixteen-bit storage, or a stated allowance;
it is not the direct result for a thirty-second eight-bit record.

The mean sample spacing under constant descent is
$\Delta s=v_d/f_c=2.1\,\mathrm m$; a $20\,\mathrm m$ layer lasts
$1.9048\,\mathrm s$ and contains approximately $9.52$ sampling intervals.
The exact integer frame count depends on the alignment of triggers and
layer boundaries. Averaging $n$ correlated errors with covariance $C$
gives
\begin{equation}
 \operatorname{Var}(\overline\epsilon)=\frac{\boldsymbol1^{\mathsf T}C\boldsymbol1}{n^2}.
\end{equation}
Only independent equal-variance errors reduce to $\sigma^2/n$; common
bias and suspension-model error do not disappear by layer averaging.

A serialized processor must satisfy $f_c\,\mathbb E[t_{\rm service}]<1$
to avoid persistent queue growth, with margin for variability and other
tasks. At $5\,\mathrm{Hz}$ the frame interval is $200\,\mathrm{ms}$;
$250\,\mathrm{ms}$ per frame is too slow for one sequential worker.
Several workers or pipelining may meet throughput, but require measured
resource contention, queue limits and timing. Distinguish latency from
throughput. For a logging stall of duration $T_{\mathrm{stall}}$, a sufficient
image buffer under constant input rate is $R_{\rm image}T_{\mathrm{stall}}$,
in addition to data already buffered; recovery requires write rate above
the mean input rate.

\subsection{Battery energy, peak current and mechanical loading}

For mission phases $j$ of duration $t_j$ with load power $P_j$ and
conversion efficiency $\eta_j$, energy conservation requires
\begin{equation}
 E_{\rm batt,required}=\sum_j\frac{P_jt_j}{\eta_j},\qquad
 E_{\rm usable}=V_{\rm nominal}Q_{\rm Ah}f_{\rm usable}
 \quad\text{in Wh if $t_j$ is in hours}.
\end{equation}
Hence
\begin{equation}
 Q_{\rm Ah}\geq\frac{\sum_j P_jt_j/\eta_j}
                    {V_{\rm nominal}f_{\rm usable}}.
\end{equation}
$f_{\rm usable}$ includes justified capacity, temperature, ageing and
cutoff effects rather than an invented universal derating. A safety or
design margin can be applied explicitly. Ground waiting and brief camera
operation must be represented by different phases. No numerical capacity
or endurance can be certified without measured power and battery data.
Peak operation also needs
\begin{equation}
 I_{\rm batt,peak}\simeq\frac{P_{\rm peak}}{\eta V_{\rm batt}},\qquad
 V_{\rm loaded}=V_{\rm open}-I_{\rm batt}R_{\rm internal},
\end{equation}
and a demonstrated minimum rail voltage during radio, compute and storage
peaks. A three-hour bench run tests the configured article, not every
temperature or battery condition.

For a specified equivalent acceleration $a_{\rm design}$, a mounted item
of mass $m_i$ applies inertial load $F_i=m_i a_{\rm design}$. If the
proposal's $20g$ is the adopted load case, $a_{\rm design}=20g$; whether
weight is added depends on the stated test convention. Stress and bending
checks require geometry, material allowables and load direction. A stated
connection pull load does not follow from terminal canopy force alone:
deployment loads require an inflation/deceleration model or measurement.
For stopping speed change $\Delta v$ over $\Delta t$, impulse balance
gives average net force $m\Delta v/\Delta t$, but does not bound peak force.
The proposal's load and endurance numbers are inputs awaiting verification
against the applicable competition rules and physical tests.

\subsection{Radio airtime with explicit parameters}

Payload byte count and packet rate alone do not establish airtime. For
the Semtech SX1276 LoRa packet convention, let $\mathrm{SF}$ denote
spreading factor, $\mathrm{BW}$ bandwidth in hertz, $P_L$ payload bytes,
$n_{\mathrm{pre}}$ programmed preamble symbols, $C$ CRC flag, $I$ implicit-header
flag, $D$ low-data-rate optimization flag, and $r_c=1,\ldots,4$ the coding
rate parameter corresponding to $4/5,\ldots,4/8$.
\begin{align}
 T_{\rm sym}&=2^{\mathrm{SF}}/\mathrm{BW},\\
 n_{\rm payload}&=8+\max\left\{
 \left\lceil\frac{8P_L-4\mathrm{SF}+28+16C-20I}
 {4(\mathrm{SF}-2D)}\right\rceil(r_c+4),\,0\right\},\\
 T_{\rm packet}&=(n_{\mathrm{pre}}+4.25+n_{\rm payload})T_{\rm sym}.
\end{align}
This is protocol accounting, not a modulation law derived from the CanSat
mechanics. For an illustrative explicit-header packet with $P_L=40$,
$\mathrm{SF}=7$, $\mathrm{BW}=125000$, $r_c=1$, $C=1$, $I=D=0$ and
$n_{\mathrm{pre}}=8$, the numerator is $336$, the ceiling is $12$,
$n_{\rm payload}=68$, and $T_{\rm packet}=82.176\,\mathrm{ms}$.
At one packet per second that example occupies $8.2176\%$ of transmission
time before retries or acknowledgements. It is not the selected flight
configuration or a compliance verdict. Account for all transmissions in
the applicable assessment interval and verify current regional sub-band
conditions for the actual settings.

The formula and register meanings are verified against the manufacturer's
datasheet, section 4.1.1.7 (Semtech document hosted by SparkFun):
\url{https://cdn.sparkfun.com/assets/7/7/3/2/2/SX1276_Datasheet.pdf}.
For link budgeting, a separate conditional calculation is
\begin{equation}
 P_{r,\rm dBm}=P_{t,\rm dBm}+G_t+G_r-L_{\rm path}-L_{\rm other},
 \quad L_{\rm FS}=20\log_{10}(4\pi R/\lambda),\quad \lambda=c/f.
\end{equation}
Free-space loss assumes an unobstructed far-field path. Can orientation,
antenna pattern, ground reflections, interference and receiver sensitivity
must be measured or bounded before claiming range.

\subsection{Numerical checks completed and remaining physical evidence}

Independent in-memory calculations performed during this review checked
the following identities. They used no flight measurements.
The calculations are now preserved as reproducible standard-library
unit tests in \texttt{sim/maths/test\_physics\_extra.py}; all ten tests
passed on execution. They additionally check the radio example, neutral
MOST limit, lower boundary, dry-adiabatic invariant, shear solution,
pendulum counterexample, storage and blur arithmetic, and source-file
control bytes.
\begin{center}
\begin{tabular}{p{0.31\linewidth}p{0.27\linewidth}p{0.30\linewidth}}
\toprule
Check & Method & Result\\
\midrule
First-order response, $r=0.5,1,2$ & RK4, $\Delta t=0.001\,\mathrm s$
 & max. error $<1.9\times10^{-13}\,\mathrm{m/s}$\\
Isothermal height, $300\,\mathrm m$ & pressure forward/inverse
 & error $<6\times10^{-13}\,\mathrm m$\\
MOST integrated correction & quadrature in log height
 & max. residual $1.5\times10^{-14}$\\
MOST gradient correction & central difference
 & max. residual $1.4\times10^{-9}$\\
\bottomrule
\end{tabular}
\end{center}
For the lag check the independently integrated equation was
$\dot u=(5-u)/\tau_h$, $u(0)=0$, for five seconds; the analytic comparator
was $5(1-e^{-t/\tau_h})$. For the hydrostatic check,
$R_d=287.05\,\mathrm{J/(kg\,K)}$, $T=288.15\,\mathrm K$ and
$p(0)=101325\,\mathrm{Pa}$. MOST checks used $\zeta=-2,-0.1,0.1,2$,
$z_0=0.1\,\mathrm m$, $|L_O|=30\,\mathrm m$, central-difference step
$10^{-6}$, and composite Simpson quadrature with $10000$ intervals.
These residuals verify algebra and numerical implementation for the chosen
models; they do not establish agreement with the real atmosphere.

Physical validation still needs time-aligned drops measuring deployment,
vertical speed, lateral response, line geometry and sensor offsets;
temperature response tests in representative airflow; suitable independent
wind observations; measured camera exposure and image signal quality;
and endurance, storage-stall and radio tests in flight configuration.
Retain measured data and report fitted coefficients and their uncertainty.
The deterministic counterexamples above are sufficient to reject
unconditional claims about horizontal lag, gyro-derived static tilt,
vertical sampling at all separations, bulk Richardson thresholds, or
Obukhov length inferred from temperature lapse alone.
